% ============================================================
% DOCUMENTCLASS
% ============================================================
\documentclass[12pt,oneside]{book}

% ============================================================
% METADATOS
% ============================================================
\newcommand{\universidad}{Universidad de Buenos Aires (UBA)}
\newcommand{\facultad}{Facultad de Derecho}
\newcommand{\carrera}{Carrera de Abogacía}
\newcommand{\materia}{Práctica Profesional de Abogacía}
\newcommand{\titulo}{El juicio de desalojo por cambio de destino}
\newcommand{\autor}{Isaac Edgar Camacho Ocampo}
\newcommand{\anio}{2026}

% ============================================================
% PAQUETES
% ============================================================
\usepackage[utf8]{inputenc}
\usepackage[T1]{fontenc}
\usepackage[spanish,es-nodecimaldot]{babel}

\usepackage[
    top=2.5cm,
    bottom=2.5cm,
    left=3cm,
    right=2.5cm,
    headheight=15pt
]{geometry}

\usepackage{amsmath}
\usepackage{amssymb}
\usepackage{mathtools}

\usepackage{graphicx}
\usepackage{float}
\usepackage{caption}
\usepackage{subcaption}

\usepackage{booktabs}
\usepackage{array}
\usepackage{longtable}
\usepackage{tabularx}

\usepackage{enumitem}

\usepackage{xcolor}
\usepackage[most]{tcolorbox}

\usepackage{fancyhdr}
\usepackage{ragged2e}

\usepackage[
    colorlinks=true,
    linkcolor=blue!50!black,
    citecolor=green!40!black,
    urlcolor=blue!60!black,
    pdftitle={\titulo},
    pdfauthor={\autor},
    pdfsubject={Apuntes universitarios},
    bookmarks=true
]{hyperref}

% ============================================================
% CONFIGURACIÓN
% ============================================================
\graphicspath{
    {./img/}
    {./graficos/}
    {./figuras/}
}

\setcounter{tocdepth}{2}

\setlength{\parindent}{0pt}
\setlength{\parskip}{0.5em}

\setlist[itemize]{
    topsep=0.3em,
    itemsep=0.2em
}

\setlist[enumerate]{
    topsep=0.3em,
    itemsep=0.2em
}

\clubpenalty=10000
\widowpenalty=10000

% ============================================================
% COMANDOS
% ============================================================
\newcommand{\abs}[1]{\left|#1\right|}
\newcommand{\norm}[1]{\left\lVert#1\right\rVert}
\newcommand{\vect}[1]{\mathbf{#1}}
\newcommand{\deriv}[2]{\frac{d#1}{d#2}}
\newcommand{\pderiv}[2]{\frac{\partial#1}{\partial#2}}

% ============================================================
% ENTORNOS / CAJAS
% ============================================================
\newtcolorbox{definition}[1]{
    enhanced,
    breakable,
    title={Definición: #1},
    fonttitle=\bfseries,
    colback=blue!3,
    colframe=blue!50!black,
    boxrule=0.8pt,
    arc=2mm
}

\newtcolorbox{important}{
    enhanced,
    breakable,
    title={Importante},
    fonttitle=\bfseries,
    colback=yellow!8,
    colframe=orange!70!black,
    boxrule=0.8pt,
    arc=2mm
}

\newtcolorbox{example}[1]{
    enhanced,
    breakable,
    title={Ejemplo: #1},
    fonttitle=\bfseries,
    colback=green!4,
    colframe=green!40!black,
    boxrule=0.8pt,
    arc=2mm
}

\newtcolorbox{formula}[1]{
    enhanced,
    breakable,
    title={Fórmula / Regla: #1},
    fonttitle=\bfseries,
    colback=purple!4,
    colframe=purple!50!black,
    boxrule=0.8pt,
    arc=2mm
}

\newtcolorbox{exercise}[1]{
    enhanced,
    breakable,
    title={Ejercicio: #1},
    fonttitle=\bfseries,
    colback=gray!5,
    colframe=gray!60!black,
    boxrule=0.8pt,
    arc=2mm
}

\newtcolorbox{note}{
    enhanced,
    breakable,
    title={Nota},
    fonttitle=\bfseries,
    colback=white,
    colframe=gray!50,
    boxrule=0.6pt,
    arc=2mm
}

% ============================================================
% CONFIGURACIÓN FINAL (encabezados y pies)
% ============================================================
\pagestyle{fancy}
\fancyhf{}

\fancyhead[L]{\nouppercase{\leftmark}}
\fancyhead[R]{\materia}

\fancyfoot[C]{\thepage}

\renewcommand{\headrulewidth}{0.4pt}
\renewcommand{\footrulewidth}{0pt}

\fancypagestyle{plain}{
    \fancyhf{}
    \fancyfoot[C]{\thepage}
    \renewcommand{\headrulewidth}{0pt}
}

% ============================================================
% DOCUMENT
% ============================================================
\begin{document}

% ---------------- PORTADA ----------------
\begin{titlepage}

\thispagestyle{empty}

\begin{center}

\vspace*{1cm}

{\large \textsc{\universidad}}

\vspace{0.2cm}

{\large \textsc{\facultad}}

\vspace{0.2cm}

{\large \carrera}

\vspace{3cm}

{\Huge \textbf{\materia}}

\vspace{1cm}

{\LARGE \titulo}

\vspace{0.8cm}

\textit{Comprender el proceso es tan necesario como conocer el derecho: juntos convierten un reclamo justo en una sentencia.}

\vspace{1.2cm}

\rule{0.70\textwidth}{0.5pt}

\vspace{0.5cm}

{\large Apuntes de estudio}

\vspace{0.5cm}

\rule{0.70\textwidth}{0.5pt}

\vfill

{\large \autor}

\vspace{0.3cm}

{\large \anio}

\end{center}

\vfill

\begin{flushleft}
\textit{Este es un modesto aporte a los estudiantes de ``Práctica Profesional de Abogacía'' no pretende sustituir el material oficial de la materia.}
\end{flushleft}

\end{titlepage}

% ---------------- ÍNDICE ----------------
\tableofcontents
\clearpage

% ============================================================
% INTRODUCCIÓN
% ============================================================
\chapter*{Introducción a la unidad temática}
\addcontentsline{toc}{chapter}{Introducción a la unidad temática}

Esta unidad reconstruye, a través de una simulación de consulta en un estudio jurídico, un caso real de conflicto locativo: una familia propietaria de un inmueble heredado busca recuperarlo porque el inquilino cambió el destino pactado —vivienda familiar— por un emprendimiento comercial no autorizado. A partir de ese caso se desarrolla el fundamento jurídico del desalojo por cambio de destino y, en su segunda mitad, el trámite práctico completo para presentar una demanda de desalojo: la estructura de la demanda, la firma hológrafa y el poder, la legitimación y la notificación, la prueba, los planteos de inconstitucionalidad y reserva del caso federal, y la presentación electrónica del expediente tanto ante el Poder Judicial de la Nación como ante la Provincia de Buenos Aires.

El contenido resulta directamente aplicable al ejercicio profesional: reproduce, paso a paso, la primera entrevista con un cliente y la primera demanda que la mayoría de los abogados jóvenes redacta en su vida profesional. Aprender a escuchar un relato familiar desordenado y traducirlo en hechos jurídicamente relevantes, a distinguir qué prueba sirve y cuál no, y a manejar los sistemas de gestión electrónica de los que hoy depende cualquier litigio, constituye una competencia de uso cotidiano en un estudio jurídico, tanto en materia de locaciones como en cualquier otro fuero. Más allá de lo profesional, resulta también de utilidad personal: comprender cómo funciona un contrato de locación, qué puede y qué no puede hacerse con un inmueble alquilado, y qué pasos previos existen antes de llegar a un juicio, es información relevante para cualquier persona, propietaria o inquilina.

\begin{important}
\textbf{Cómo se relacionan los bloques del contenido.} Puede pensarse la estructura de esta unidad como la construcción de una casa. Primero es necesario conocer el terreno y a la familia que la habita —los hechos del caso—; después hay que verificar que existan los cimientos jurídicos que sostienen el reclamo —el fundamento del desalojo por cambio de destino—; luego se levantan las paredes —la demanda y el expediente electrónico—; se refuerza la estructura con la prueba y los planteos procesales; y, por último, se llega a la etapa en la que el expediente efectivamente avanza en el sistema del tribunal. Cada bloque es una etapa necesaria de un mismo proceso: sin hechos claros no hay demanda posible, sin demanda bien armada no hay prueba que valga, y sin conocer el trámite electrónico ni la demanda mejor redactada llega a destino.
\end{important}

% TODO: contenido incompleto o ambiguo en las notas originales — la nota metodológica original advierte que el audio presenta degradaciones propias de una transcripción automática (cortes de palabras, nombres mal reconocidos y pasajes ininteligibles); el presente documento reconstruye fielmente el contenido conceptual, los ejemplos, las normas citadas y el desarrollo cronológico de la clase, conservando en cursiva las expresiones textuales de la docente cuando se citan.

% ============================================================
% CAPÍTULO 1
% ============================================================
\chapter{La consulta: el caso y los hechos}

\section{La simulación como método de enseñanza}

La unidad se desarrolla mediante una actividad de simulación profesional: un grupo de estudiantes conforma un ``bufete'' o estudio jurídico y representa una consulta de clientes ante una abogada, siguiendo un guion previamente escrito. El personaje de la abogada es interpretado por la docente a cargo, mientras que el resto de los participantes representa a los distintos clientes.

\begin{note}
La elección del formato de role-play, en lugar de una exposición lineal del caso, responde a una decisión pedagógica deliberada: presentar la conflictiva de una manera que permita observar cómo llega la información desde el cliente, con sus múltiples intereses entrecruzados, antes de traducirla en categorías jurídicas.
\end{note}

Antes de dar comienzo a la escena, se advierte a los estudiantes que los nombres y la ambientación del caso —vinculados a distintos clubes de fútbol— fueron elegidos de manera intencional como recurso didáctico para hacer más ameno el relato de los hechos, sin que ello disminuya la relevancia jurídica del caso planteado.

\section{Los hechos del caso: el inmueble, el contrato y el conflicto}

\subsection{El origen del inmueble y los sujetos intervinientes}

La escena simulada representa la primera entrevista de un estudio jurídico con una familia que consulta por un conflicto vinculado a un inmueble dado en locación. Concurren tres clientes: Marta Riquel (madre, viuda) y sus dos hijos, Juan Román Riquel y Claudia Riquel, quienes se presentan como copropietarios del bien en carácter de herederos.

El inmueble había pertenecido a los abuelos paternos de Juan Román y Claudia, quienes lo adquirieron con el fruto de su trabajo tras emigrar desde Italia. Tras el fallecimiento del esposo de Marta, el bien pasó a integrar el acervo hereditario de los tres consultantes, encontrándose la sucesión ya iniciada al momento de la consulta.

\begin{important}
\textbf{Dos conflictos, una sola consulta.} La entrevista con el cliente rara vez trae un único problema jurídico. En este caso conviven una cuestión sucesoria (la disputa entre los coherederos sobre vender o conservar el bien) y una cuestión contractual y procesal (recuperar la tenencia frente al locatario). Distinguir ambos planos desde el primer contacto —sin descartar ninguno, pero priorizando el que requiere acción inmediata— es una destreza central de la entrevista de admisión.
\end{important}

Dentro de la familia coexisten posiciones distintas sobre el destino final del inmueble: Claudia y su madre pretenden conservar la propiedad por su valor sentimental, mientras que Juan Román necesita hacerse de su parte indivisa para afrontar deudas personales. La abogada identifica que se trata de dos conflictos jurídicamente distintos que conviven en la misma consulta —la disputa sucesoria entre los tres herederos y la situación con el locatario que ocupa la casa— y decide priorizar el segundo, sin descuidar el primero, atendiendo a que se trata del reclamo que requiere acción inmediata.

\subsection{El contrato de locación y la cláusula de destino}

El inmueble se encontraba alquilado a Martín Milán y a quien se identifica como su media hermana. El contrato de locación se encontraba vigente al momento de la consulta y el canon locativo se hallaba al día, extremo que resulta relevante porque descarta, desde el inicio, la posibilidad de fundar el reclamo en la falta de pago.

\begin{definition}{Cláusula de destino}
Estipulación contractual mediante la cual las partes fijan el uso específico y exclusivo que se dará al inmueble locado. En el contrato analizado, la cláusula tercera establece: \textit{``Destino: vivienda familiar exclusivamente.''} La abogada verifica personalmente esta cláusula leyendo el contrato original que le acercan los clientes, antes de avanzar con cualquier estrategia procesal.
\end{definition}

\subsection{El incumplimiento: cambio de destino del inmueble}

Pese a que el contrato se encuentra vigente y el alquiler se paga puntualmente, el inmueble dejó de destinarse a vivienda familiar. La pareja del locatario, identificada en el caso como Laura Bataglia, instaló dentro de la vivienda un emprendimiento comercial de venta de productos de un club de fútbol.

\begin{example}{El emprendimiento dentro de la vivienda familiar}
Según relatan los consultantes, existen compradores que forman filas porque el comercio no vende en línea; se venden camisetas, buzos, equipos deportivos, tazas y banderines; se utiliza la cocina, el horno y el gas de la vivienda para preparar productos alimenticios con los colores y el escudo del club; los clientes ingresan a la vivienda y también esperan afuera, entre cincuenta y cien personas los fines de semana; hay vehículos estacionados en doble fila y transporte que traslada compradores; y se escuchan música, cánticos e incluso insultos a otros equipos a cualquier hora del día. La familia cuenta con fotografías, videos y mensajes de los propios vecinos que acreditan la situación.
\end{example}

Consultado sobre el reclamo, el locatario reconoce conocer la intención de la familia de recuperar la tenencia, pero se niega a restituir el inmueble invocando que el contrato se encuentra vigente. Respecto de la actividad comercial, atribuye su desarrollo exclusivamente a su pareja, desligándose de responsabilidad personal.

\subsection{La garantía del contrato}

El contrato de locación fue garantizado con una garantía real constituida por un matrimonio, titular del inmueble ofrecido en garantía. Ambos garantes firmaron el contrato, con sus firmas certificadas por escribano público, y se había solicitado con anterioridad un informe de dominio para verificar la titularidad del bien dado en garantía y la inexistencia de gravámenes o restricciones.

\subsection{Síntesis de los elementos reunidos en la entrevista}

Al cierre de la primera entrevista, la abogada cuenta con los siguientes elementos: la historia del inmueble y su origen hereditario; la sucesión abierta entre los tres consultantes; el contrato de locación vigente, con destino exclusivo de vivienda familiar pactado en su cláusula tercera; el incumplimiento acreditable mediante fotografías, videos y testimonios de vecinos; y la garantía real debidamente constituida y certificada. A partir de estos elementos, se define como paso siguiente reunir toda la documentación disponible —contrato, intimaciones si las hubiera, fotografías, videos, mensajes de vecinos— e intentar, en la medida de lo posible, un diálogo previo con la parte locataria antes de promover la demanda.

La dirección del inmueble, dato relevante para determinar la competencia territorial del futuro juicio, corresponde al barrio de Núñez, Ciudad Autónoma de Buenos Aires.

% ============================================================
% CAPÍTULO 2
% ============================================================
\chapter{Fundamentos del desalojo por cambio de destino}

\section{El derecho del locador y la prueba del cambio de destino}

El derecho del locador a exigir el desalojo cuando el inquilino modifica el destino pactado en el contrato constituye una causal autónoma, distinta de la falta de pago y del vencimiento del plazo contractual.

\begin{formula}{Causal de desalojo por cambio de destino}
El locador tiene derecho a demandar el desalojo cuando el inquilino modifica el destino contractualmente pactado para el inmueble, con independencia de que el canon locativo se encuentre al día y de que el plazo del contrato no haya vencido. La causal a invocar en la demanda es específicamente el \textbf{desalojo por uso indebido, uso prohibido o cambio de destino}, y no la falta de pago ni el vencimiento del plazo.
\end{formula}

\begin{example}{Cómo se traduce el relato familiar en una causal jurídica}
El uso intensivo y comercial de la cocina y del suministro eléctrico, la venta a público masivo y las quejas vecinales documentadas no constituyen datos anecdóticos: son la prueba concreta de que el inmueble dejó de destinarse a vivienda familiar para destinarse a una actividad comercial. Si se cocina y se vende a gran público, el uso del inmueble deja de asimilarse al de una vivienda.
\end{example}

Como elemento probatorio adicional, puede requerirse por oficio a la dependencia de habilitaciones del gobierno local un informe sobre si la zona en la que se ubica el inmueble habilita actividad comercial o si se trata de una zona residencial. Si la actividad no resulta habilitada como zona comercial, queda reforzado el derecho a iniciar la demanda por cambio de destino.

\begin{important}
\textbf{Tres causales, una sola herramienta procesal, distinta según el objetivo.} El desalojo por falta de pago, el desalojo por vencimiento del plazo y el desalojo por cambio de destino comparten la misma vía procesal, pero requieren pruebas y planteos distintos. En el caso analizado, la causal aplicable es el cambio de destino: el contrato está vigente y el alquiler al día, pero el inmueble se utiliza para un fin distinto del pactado. Antes de demandar, siempre conviene evaluar si el cliente busca recuperar el inmueble (desalojo) o solo hacer cesar el uso indebido conservando la locación (intimación de cese).
\end{important}

Cuando el objetivo del cliente no es recuperar el inmueble sino simplemente hacer cesar la actividad comercial y mantener la locación vigente, existe una alternativa menos drástica: intimar por carta documento el cese de la actividad, invocando la cláusula contractual de destino exclusivo de vivienda familiar, bajo apercibimiento de iniciar la acción de desalojo por cambio de destino si el requerimiento no es atendido.

\section{Intimación previa, carta documento y mediación}

Corresponde distinguir en qué casos la intimación previa a la demanda resulta jurídicamente obligatoria y en qué casos constituye únicamente una buena práctica profesional.

\begin{table}[H]
\centering
\caption{Intimación previa y mediación en el desalojo}
\begin{tabularx}{\textwidth}{>{\raggedright\arraybackslash}p{0.28\textwidth} >{\raggedright\arraybackslash}p{0.26\textwidth} X}
\toprule
\textbf{Instrumento} & \textbf{¿Es obligatorio en este caso?} & \textbf{Función} \\
\midrule
Carta documento (intimación de pago) & Obligatoria en desalojo por falta de pago & Constituir en mora al locatario antes de demandar \\
Carta documento (cambio de destino) & No obligatoria; buena práctica recomendada & Acercar a las partes e intentar una solución sin litigio \\
Mediación previa & No obligatoria en juicios de desalojo & Vía alternativa de resolución de conflictos, a criterio del profesional \\
\bottomrule
\end{tabularx}
\end{table}

A diferencia de otros procesos en los que la mediación constituye un requisito de admisibilidad de la demanda, en los juicios de desalojo la mediación no es obligatoria y queda a criterio de cada profesional proponerla como vía alternativa antes de litigar.

\begin{important}
\textbf{Intimar no siempre es obligatorio, pero casi siempre conviene.} La distinción entre lo jurídicamente exigible y lo profesionalmente recomendable es una de las claves de la práctica: aunque ni la carta documento por cambio de destino ni la mediación sean requisitos legales para demandar el desalojo, ambas herramientas permiten intentar una solución más rápida y económica para el cliente, y dejan constancia documental de la buena fe procesal de quien las propone.
\end{important}

% ============================================================
% CAPÍTULO 3
% ============================================================
\chapter{La demanda y el expediente electrónico}

\section{Estructura de la demanda y límites del expediente electrónico}

El artículo 330 del Código Procesal Civil y Comercial de la Nación fija los requisitos formales mínimos de todo escrito de demanda —nombre y domicilio del demandante, nombre y domicilio del demandado, objeto, hechos y derecho, petición—, pero no constituye una plantilla suficiente para casos complejos.

\begin{important}
\textbf{Artículo 330, Código Procesal Civil y Comercial de la Nación.} No corresponde basarse únicamente en los lineamientos del artículo 330 —que resulta acotado— sino desarrollar apartados adicionales de acuerdo con la complejidad de la demanda. Un caso con varias causales y varios sujetos, como el analizado, exige rubros que ese artículo no prevé expresamente.
\end{important}

\subsection{El límite de tamaño de los archivos digitales}

Al presentar el escrito de demanda junto con la prueba documental, el sistema de gestión electrónica exige que cada archivo adjunto no supere un tamaño máximo determinado.

\begin{example}{Qué hacer cuando la prueba documental supera el límite}
Si la documentación a acompañar supera el peso máximo admitido por el sistema, la solución práctica consiste en dividir el archivo en partes —``documental uno'', ``documental dos''— y adjuntar cada parte por separado al proceso judicial, de modo que el sistema permita la carga de cada archivo individualmente.
\end{example}

\begin{note}
El límite de tamaño de archivo se fue actualizando con las sucesivas acordadas de la Corte Suprema de Justicia de la Nación: comenzó en 5 MB y hoy se encuentra en 20 MB, de acuerdo con las acordadas de gestión electrónica dictadas a partir del año 2020, en el marco de la aceleración de la migración del expediente en papel al expediente electrónico.
\end{note}

\begin{table}[H]
\centering
\caption{Fuentes normativas del expediente electrónico}
\begin{tabularx}{\textwidth}{>{\raggedright\arraybackslash}p{0.27\textwidth} >{\raggedright\arraybackslash}p{0.27\textwidth} X}
\toprule
\textbf{Fuente normativa} & \textbf{Sistema que regula} & \textbf{Vigencia} \\
\midrule
Acordadas de la Corte Suprema (desde 2020) & Gestión electrónica del expediente & Se actualizan periódicamente; el límite de archivo vigente es de 20 MB \\
Código Procesal Civil y Comercial de la Nación & Formas y plazos de los escritos judiciales (sistema originalmente analógico) & Vigente; se aplica junto con las acordadas electrónicas \\
Reglamento de la Justicia Nacional & Formas complementarias de los escritos judiciales & De origen anterior, pero aún vigente y de aplicación práctica \\
\bottomrule
\end{tabularx}
\end{table}

\subsection{La lentitud del despacho de los escritos electrónicos}

La transición del sistema analógico al electrónico trajo consigo un problema práctico relevante: la demora de los juzgados para aceptar y despachar los escritos electrónicos, que quedan a la espera de que el personal judicial los reciba desde su propio sistema informático, denominándose a esa espera ``dependencia''. Existen juzgados que despachan los escritos en plazos prudenciales y otros que acumulan demoras de hasta un mes.

\begin{example}{El seguimiento personal del escrito}
Cuando un escrito urgente no se despacha, una práctica habitual consiste en concurrir personalmente a la mesa de entradas del juzgado para solicitar, de manera directa, el despacho del escrito pendiente, dado que el sistema electrónico no siempre garantiza atención inmediata por parte del personal a cargo.
\end{example}

\begin{important}
\textbf{La norma no alcanza sin conocer la práctica del fuero.} Cumplir el artículo 330 y las acordadas de gestión electrónica es condición necesaria, pero no suficiente: conocer el límite de tamaño de archivo, saber dividir la prueba documental, y tener presente que el despacho de un escrito depende de que el personal del juzgado lo acepte desde su propio sistema, son datos de práctica profesional que no figuran en ningún código pero que determinan si un escrito llega a horario o se pierde.
\end{important}

\section{Firma hológrafa, firma electrónica y poder}

Los tribunales suelen rechazar escritos firmados electrónicamente por el propio cliente, exigiendo en su lugar la firma hológrafa —manuscrita— digitalizada.

\begin{example}{Firma hológrafa escaneada, no fotografiada}
Cuando el cliente no puede concurrir en persona a firmar, el escrito se envía a la persona designada para gestionar la firma; esa persona hace firmar de puño y letra al cliente, escanea el documento en formato PDF —nunca como fotografía— y lo reenvía al abogado responsable, quien recién entonces lo sube al expediente electrónico.
\end{example}

Junto con la firma hológrafa, resulta una práctica de organización interna relevante llevar una carpeta paralela de seguimiento del expediente, sin confiar exclusivamente en el sistema informático del Poder Judicial.

\begin{example}{El escrito que nadie había despachado}
En un caso de la práctica profesional, se sostenía de memoria que determinado escrito ya había sido resuelto por el juzgado. Al revisar el expediente con calma, se comprobó que la resolución que se creía reciente correspondía en realidad a un escrito de años atrás, y que el planteo verdaderamente pendiente nunca había sido despachado, adviertiéndose esto recién cuando el caso llegó a la Cámara de Apelaciones por otro planteo distinto. La conclusión es que la memoria puede fallar, por lo que resulta indispensable dejar todo por escrito y revisar con cuidado el estado del expediente.
\end{example}

Esta misma lógica de respaldo documental se extiende a los hechos que el propio cliente relata verbalmente durante la entrevista, que muchas veces resultan incompletos o se omiten por vergüenza o por desconocimiento de su relevancia jurídica. En cuestiones de familia o de orden público, resulta recomendable requerir al cliente un relato escrito y firmado, o bien un correo electrónico enviado por el propio cliente, en lugar de basarse únicamente en lo escuchado durante la entrevista.

\begin{important}
\textbf{Artículo 47, Código Procesal Civil y Comercial de la Nación, y reforma del Código Civil y Comercial (2015).} El Código Procesal exigía, en su artículo 47, escritura de poder general, judicial o especial. El Código Civil y Comercial, vigente desde agosto de 2015, flexibiliza ese requisito, permitiendo que el poder para el abogado se otorgue por instrumento privado o por instrumento privado con firma certificada. En cualquier caso, el poder debe acompañarse siempre en la primera presentación junto con el escrito de demanda, para que el juez pueda verificar la personería; de lo contrario, el demandado puede oponer la excepción de falta de personería.
\end{important}

En cuestiones de familia y en demandas de inicio en materia patrimonial, resulta preferible que el propio cliente firme la demanda —presentándose primero con patrocinio letrado— y que recién después se incorpore el poder para actuar como apoderado, de modo que quede constancia de la voluntad del cliente desde el inicio del proceso.

\begin{example}{El divorcio de común acuerdo que derivó en una demanda contra la propia abogada}
En un caso de la práctica profesional anterior a la reforma de 2015, se patrocinó un divorcio que se tramitaba en apariencia de común acuerdo. Tiempo después, una de las partes decidió plantear la nulidad de lo actuado y demandó también a la profesional interviniente, alegando haber sido obligada a firmar el convenio. El reclamo pudo desestimarse por dos razones: la clienta había firmado el escrito como patrocinada, habiéndolo leído ella misma, y esa firma había sido puesta ante un escribano público que certificó su autenticidad.
\end{example}

\begin{important}
\textbf{La firma hológrafa certificada protege tanto al cliente como al abogado.} Exigir una firma manuscrita, certificada o acompañada de un relato escrito del cliente no es solamente un requisito formal para que el tribunal acepte el escrito: es también una herramienta de resguardo profesional frente a reclamos posteriores del propio cliente.
\end{important}

\section{Legitimación, personería y notificación}

Corresponde distinguir con precisión dos excepciones procesales que suelen confundirse entre sí: la falta de personería y la falta de legitimación pasiva para obrar.

\begin{table}[H]
\centering
\caption{Falta de personería y falta de legitimación pasiva}
\begin{tabularx}{\textwidth}{>{\raggedright\arraybackslash}p{0.24\textwidth} >{\raggedright\arraybackslash}p{0.37\textwidth} X}
\toprule
\textbf{Excepción procesal} & \textbf{Qué cuestiona} & \textbf{Ejemplo del caso} \\
\midrule
Falta de personería & Que quien actúa en el proceso (letrado o representante) no acreditó debidamente su representación o poder & El abogado de los consultantes no acompaña el poder de dos de los coherederos junto con la demanda \\
Falta de legitimación pasiva para obrar & Que el demandado no tiene ninguna relación jurídica con el objeto del litigio & Se demanda a un amigo del locatario que solo visita la casa, pero no la ocupa ni la usa \\
\bottomrule
\end{tabularx}
\end{table}

En un juicio de desalojo corresponde demandar no solo al locatario formal, sino también a los subinquilinos y a todo ocupante real del inmueble, aunque no figure como parte firmante del contrato.

\begin{example}{El croquis como prueba de contexto}
Resulta útil acompañar a la demanda un croquis, un enlace a un mapa satelital y fotografías del entorno del inmueble, de modo que el juzgado pueda verificar que se trata de una zona exclusivamente residencial y no de una arteria comercial. Aunque el contrato ya especifique que el destino pactado es de vivienda, estos elementos adicionales refuerzan la convicción del juzgado, sin ser en sí mismos prueba imprescindible.
\end{example}

Respecto de la notificación del traslado de la demanda, cuando el oficial notificador no logra notificar al demandado en el domicilio denunciado —por ejemplo, porque se niega su presencia o la numeración no coincide—, corresponde acreditar que la persona efectivamente reside allí y solicitar que se notifique bajo responsabilidad de la parte actora. Bajo esa modalidad, si nuevamente se niega la residencia del demandado, el oficial notificador fija la cédula en la puerta de acceso del inmueble.

\begin{example}{El croquis para el oficial notificador}
Para reducir el margen de error del oficial notificador, resulta recomendable adjuntar a la cédula fotografías y un croquis detallado que identifique con precisión el domicilio a notificar, incluyendo referencias visuales concretas del entorno inmediato del inmueble.
\end{example}

\begin{important}
\textbf{La precisión material como estrategia procesal.} Un croquis detallado, con referencias visuales concretas del entorno, no es un capricho estético: reduce el margen de error del oficial notificador, evita devoluciones de cédula por ``domicilio no ubicado'' y acelera el trámite, ahorrando meses en un proceso donde cada notificación fallida implica volver a empezar el trámite de diligenciamiento.
\end{important}

% ============================================================
% CAPÍTULO 4
% ============================================================
\chapter{La prueba y los planteos procesales}

\section{Fuentes y medios de prueba}

La demanda constituye un acto procesal de postulación, categoría que se distingue de los actos procesales de notificación —a cargo del oficial notificador— y de los actos procesales de obtención, mediante los cuales se solicita un informe.

Antes de tipear directamente la demanda existe una etapa previa de diligencias preliminares, reguladas en el artículo 323 del Código Procesal Civil y Comercial de la Nación, que habilita un conjunto de medidas para asegurar datos o documentos necesarios para encuadrar correctamente la futura demanda.

Como premisa general de la práctica litigiosa, la parte demandada suele negar, rechazar y desconocer la totalidad de los hechos y de la prueba documental, incluso frente a documentación certificada por escribano público, con el objeto de dilatar el proceso solicitando, por ejemplo, prueba pericial caligráfica improcedente cuando la firma ya se encuentra certificada por funcionario público.

\begin{definition}{Fuente de prueba}
Lugar donde queda registrado el hecho que se quiere probar. Es anterior y externa al proceso. Por ejemplo, la percepción directa de un vecino respecto de lo que observó u oyó.
\end{definition}

\begin{definition}{Medio de prueba}
Mecanismo que el Código habilita para incorporar una fuente de prueba al expediente judicial —documental, testimonial, informativa, informática—. Se produce dentro del proceso, según lo que el Código regula.
\end{definition}

\begin{example}{El vecino como fuente, el testigo como medio}
Si existen videos y fotografías, esas constituyen fuentes ya registradas. Pero si se pretende incorporar al proceso lo que percibieron los vecinos —los olores, la cantidad de gente, los cánticos—, no corresponde librar un oficio al vecino: corresponde ofrecer prueba testimonial, porque la fuente inmediata es la percepción del propio vecino, y el medio que el Código prevé para incorporarla al proceso es el testimonio.
\end{example}

\begin{important}
\textbf{El derecho de fondo no alcanza sin dominio del proceso.} Tener razón en el fondo del asunto —que el contrato prohíbe la actividad comercial, que hay quejas documentadas de los vecinos— no garantiza ganar el juicio si la prueba se ofrece por el medio equivocado. Si la fuente y el medio de prueba están mal pedidos, esa prueba no sirve. Para litigar bien cualquier proceso especial como el desalojo es indispensable dominar antes los institutos generales del proceso civil.
\end{important}

\section{Planteo de inconstitucionalidad y reserva del caso federal}

El control de constitucionalidad en el sistema argentino es difuso: no es necesario esperar a que la causa llegue a la Corte Suprema para plantear la inconstitucionalidad de una norma, sino que puede formularse ya ante el juez de primera instancia.

\begin{important}
\textbf{Oportunidad procesal del planteo de inconstitucionalidad.} El planteo de inconstitucionalidad debe formularse en la primera presentación procesal —típicamente, al contestar demanda—. Si no se plantea en ese momento, el planteo deviene extemporáneo y ya no puede introducirse válidamente en etapas posteriores del proceso. Una vez formulado, debe mantenerse expresamente en todos los escritos relevantes posteriores para no perder la vía de eventual recurso extraordinario ante la Corte Suprema, lo que constituye la denominada reserva del caso federal.
\end{important}

Corresponde señalar un criterio propio de esta cátedra respecto de una práctica extendida en el foro: la inclusión automática, en cualquier escrito, de la cláusula de estilo ``se hace reserva del caso federal'', sin que exista una lesión constitucional real que la justifique. Según este criterio, los derechos no se reservan: cuando están vulnerados, se ejercen para obtener su protección efectiva. Reservar un eventual daño para un posible caso federal sin sustento concreto constituye una fórmula vacía, copiada de otros escritos sin verificar su procedencia. Conforme a esta postura, la reserva del caso federal solo corresponde formularse y sostenerse cuando efectivamente se está lesionando una garantía constitucional y existe un planteo de inconstitucionalidad en curso.

\begin{important}
\textbf{No toda cláusula de estilo agrega valor al escrito.} La reserva del caso federal, cuando se incluye de manera automática y sin sustento en una lesión constitucional concreta, no protege ningún derecho: es una fórmula vacía copiada de otros escritos. Corresponde reservarla —y sostenerla en cada instancia— únicamente cuando existe efectivamente un planteo de inconstitucionalidad en danza, formulado en tiempo y forma desde la primera presentación.
\end{important}

% ============================================================
% CAPÍTULO 5
% ============================================================
\chapter{El trámite ante los tribunales}

\section{La tasa de justicia}

Existen dos vías disponibles para el pago de la tasa de justicia: el formulario tradicional en papel, que se abona en una sucursal bancaria, y el pago electrónico mediante el sistema de banca judicial, para el cual el cliente debe contar con una tarjeta de débito propia o de un familiar.

\begin{example}{La resistencia del cliente a pagar con tarjeta}
Buena parte de los clientes prefiere todavía el trámite presencial con el formulario de papel, por desconfianza a ingresar los datos de su tarjeta de débito en el sistema electrónico del banco. En esos casos, el estudio entrega el formulario para que el cliente concurra a cualquier sucursal bancaria habilitada, sin necesidad de acudir a una sucursal específica.
\end{example}

\begin{important}
\textbf{Ofrecer siempre las dos opciones de pago.} El pago electrónico de la tasa de justicia es más rápido, pero no todos los clientes están dispuestos a usarlo. Conviene ofrecer ambas alternativas desde el principio, en lugar de imponer una sola vía, anticipando al cliente los inconvenientes prácticos que puede encontrar en cada una.
\end{important}

\section{La presentación electrónica ante el Poder Judicial de la Nación}

El inicio de una causa ante el Poder Judicial de la Nación (PJN) comienza con el envío de un formulario de sorteo a la Cámara del fuero correspondiente, que asigna aleatoriamente el juzgado y el número de expediente. El número asignado puede conocerse por dos vías: la respuesta automática enviada por correo electrónico a la casilla desde la que se remitió el formulario de sorteo, o la consulta directa del apartado de radicaciones del sistema del PJN, donde figura el listado de las últimas causas sorteadas.

\begin{important}
\textbf{Artículos 14 y 17, Código Procesal Civil y Comercial de la Nación.} El sorteo del juzgado responde a que la distribución de la competencia es una cuestión de orden público, ajena a la voluntad de las partes: no puede elegirse a qué juez o jueza corresponderá la causa. No obstante, el Código prevé herramientas para apartar a un magistrado del conocimiento de la causa cuando existan motivos fundados. La recusación sin causa (artículo 14) se encuentra limitada en los procesos de desalojo, mientras que la recusación con causa (artículo 17) permanece siempre disponible ante motivos objetivos, como el parentesco entre el juez y una de las partes.
\end{important}

Una vez conocido el número de expediente, la carga de la demanda y de su documentación se realiza dentro del sistema, ingresando a la sección de escritos nuevos, seleccionando la jurisdicción correspondiente e identificando el expediente por su número, no por el nombre de las partes.

\begin{formula}{Formato de carga documental en el PJN}
Toda la documentación debe subirse en formato PDF; no se admite formato Word ni imágenes en formato JPG. Cualquier documento debe convertirse previamente a PDF, respetando el límite de tamaño de archivo vigente (20 MB), dividiéndolo en partes si resulta necesario.
\end{formula}

Para el seguimiento del expediente, el sistema del PJN permite localizar la causa por el apellido de la parte actora en el apartado de consulta de causas, donde se identifican los distintos objetos posibles dentro de un juicio de desalojo —desalojo por falta de pago, desalojo por vencimiento de contrato, desalojo por otras causales, categoría en la que se encuadra el cambio de destino—. Al tratarse de un sistema de preclusión, resulta indispensable un seguimiento periódico de cada causa activa para no perder los plazos procesales.

Adicionalmente, el sistema habilita, únicamente los días martes y viernes, la opción de ``dejar nota'' respecto de las resoluciones no notificadas por otra vía, mecanismo de notificación electrónica que reemplaza a la antigua nota presencial en los estudios de tribunales y que permite evitar considerarse notificado por ministerio de la ley respecto de una resolución determinada.

\begin{important}
\textbf{Sorteo, carga y seguimiento: tres momentos que no se pueden saltear.} El circuito completo ante el PJN exige respetar tres momentos sucesivos: primero, aceptar el sorteo del juzgado como una asignación de orden público, salvo causal de recusación; segundo, cargar la demanda y toda la prueba documental en PDF dentro del expediente correcto, identificado por número y no por nombre; y tercero, hacer un seguimiento activo y periódico de las resoluciones, porque el sistema de preclusión no perdona la inacción, y los días de nota electrónica —martes y viernes— determinan cuándo se considera notificada una parte que no accedió antes a la resolución.
\end{important}

\section{La presentación electrónica ante la Provincia de Buenos Aires}

El sistema de la Provincia de Buenos Aires presenta diferencias relevantes respecto del sistema nacional. El acceso se realiza mediante token para quienes cuentan con domicilio digital habilitado, o a través de la Mesa de Entradas Virtual de la Suprema Corte de la Provincia de Buenos Aires.

A diferencia del sistema nacional, en el sistema provincial no se envía un formulario de sorteo previo: la causa se inicia completando directamente cada campo del formulario electrónico, incluyendo el departamento judicial en el que tramitará, cuya elección depende de la competencia territorial y, eventualmente, de un pacto de foro de prórroga incluido en el propio contrato de locación.

\begin{table}[H]
\centering
\caption{Diferencias entre el sistema nacional y el sistema provincial}
\begin{tabularx}{\textwidth}{>{\raggedright\arraybackslash}p{0.22\textwidth} >{\raggedright\arraybackslash}p{0.37\textwidth} X}
\toprule
\textbf{Aspecto} & \textbf{Poder Judicial de la Nación} & \textbf{Provincia de Buenos Aires (SCBA)} \\
\midrule
Inicio de la causa & Formulario de sorteo enviado previamente a la Cámara del fuero & Carga directa de todos los campos e ítems, sin sorteo previo \\
Texto de la demanda & Se adjunta como archivo PDF & Se copia y pega directamente en un campo del formulario; solo la versión firmada hológrafamente se adjunta como archivo \\
Documentación probatoria & Se agrupa en archivos de hasta 20 MB, divididos si es necesario & Un archivo en PDF por cada documento (contrato, informe de dominio, fotografías, por separado) \\
Firma y envío & Requiere firma hológrafa escaneada; el trámite pasa por la aceptación del juzgado & Firma digital registrada; el envío es directo, sin intermediación de correo ni radicaciones \\
\bottomrule
\end{tabularx}
\end{table}

El sistema provincial elimina varios pasos manuales propios del sistema nacional: no es necesario enviar un correo electrónico previo, esperar su respuesta, consultar el apartado de radicaciones, ni recordar subir la demanda una vez sorteada la causa —cuya omisión, en el sistema nacional, provoca la caducidad del sorteo y obliga a sortear nuevamente—.

\begin{important}
\textbf{Dos sistemas, una misma exigencia de rigor profesional.} Más allá de sus diferencias de diseño —sorteo previo o carga directa, archivo único o archivos separados, firma hológrafa escaneada o firma digital registrada—, tanto el sistema nacional como el provincial exigen del abogado litigante el mismo nivel de atención: conocer el circuito exacto de cada fuero, respetar los formatos exigidos y hacer un seguimiento activo del expediente hasta que la resolución llegue a destino.
\end{important}

% ============================================================
% CUADRO CORNELL
% ============================================================
\chapter{Cuadro Cornell de repaso general}

El siguiente cuadro reúne, en el formato del método Cornell, las preguntas centrales que permiten repasar activamente el contenido desarrollado a lo largo de la unidad temática.

\begin{longtable}{
    >{\raggedright\arraybackslash}p{0.30\textwidth}
    >{\raggedright\arraybackslash}p{0.60\textwidth}
}
\toprule
\textbf{Preguntas / palabras clave} & \textbf{Notas / respuestas} \\
\midrule
\endfirsthead
\toprule
\textbf{Preguntas / palabras clave} & \textbf{Notas / respuestas} \\
\midrule
\endhead

\textbf{¿Cuándo procede el desalojo por cambio de destino?} &
Cuando el locatario utiliza el inmueble para un fin distinto del pactado en el contrato —en el caso, vivienda familiar exclusiva—, aunque el alquiler esté al día y el plazo contractual no haya vencido. Es una causal autónoma, distinta de la falta de pago y del vencimiento del contrato. \\

\textbf{¿Es obligatoria la intimación previa por carta documento?} &
Solo en el desalojo por falta de pago. En el desalojo por cambio de destino no es obligatoria, pero constituye una buena práctica profesional para intentar acercar a las partes antes de litigar. \\

\textbf{¿Es obligatoria la mediación previa en los juicios de desalojo?} &
No. A diferencia de otros procesos, la mediación previa es optativa en el desalojo; queda a criterio del profesional proponerla como vía alternativa. \\

\textbf{¿Alcanza el artículo 330 CPCCN para redactar cualquier demanda?} &
No. El artículo 330 fija los requisitos mínimos, pero casos complejos, con varias causales y varios sujetos, exigen desarrollar rubros adicionales que ese artículo no prevé expresamente. \\

\textbf{¿Qué límite de tamaño tienen los archivos que se suben al expediente electrónico del PJN?} &
Veinte megabytes (20 MB) por archivo, de acuerdo con las acordadas de gestión electrónica de la Corte Suprema (el límite era de 5 MB al inicio del sistema). Si la documental supera ese peso, se divide en archivos separados (documental uno, documental dos). \\

\textbf{¿Por qué los tribunales rechazan escritos firmados con firma electrónica del cliente?} &
Porque exigen firma hológrafa (manuscrita) del cliente, escaneada en PDF —nunca fotografiada—, incluso cuando existe firma electrónica disponible. El estudio gestiona esa firma en papel y la reenvía digitalizada. \\

\textbf{¿Qué forma exige hoy el poder para el abogado?} &
Antes del Código Civil y Comercial de 2015, el artículo 47 CPCCN exigía escritura de poder. Desde agosto de 2015, el poder puede otorgarse por instrumento privado o por instrumento privado con firma certificada, pero siempre debe acompañarse en la primera presentación. \\

\textbf{¿Qué diferencia hay entre falta de personería y falta de legitimación pasiva?} &
La falta de personería cuestiona que quien actúa no acreditó su representación o poder. La falta de legitimación pasiva cuestiona que el demandado no tiene ninguna relación jurídica real con el objeto del litigio (por ejemplo, un amigo del locatario que no ocupa el inmueble). \\

\textbf{¿A quién hay que demandar en un juicio de desalojo?} &
No solo al locatario que figura en el contrato, sino también a los subinquilinos y a todo ocupante real del inmueble, aunque no hayan firmado el contrato de locación. \\

\textbf{¿Qué diferencia hay entre fuente y medio de prueba?} &
La fuente es el lugar donde quedó registrado el hecho (por ejemplo, la percepción de un vecino). El medio es el mecanismo que el Código habilita para incorporar esa fuente al proceso (documental, testimonial, informativa, informática). Si se pide el medio equivocado, la prueba no sirve. \\

\textbf{¿Cuándo corresponde plantear la inconstitucionalidad de una norma?} &
En cualquier instancia, porque el control de constitucionalidad en Argentina es difuso: puede plantearse ya ante el juez de primera instancia. Debe formularse en la primera presentación, bajo pena de resultar extemporáneo, y sostenerse en los escritos posteriores. \\

\textbf{¿Cuándo corresponde reservar el caso federal?} &
Únicamente cuando existe una lesión constitucional real y un planteo de inconstitucionalidad efectivamente en curso. Incluirla como cláusula de estilo automática, sin sustento, es una práctica criticada por la cátedra. \\

\textbf{¿Cómo se inicia una causa ante el Poder Judicial de la Nación?} &
Se envía un formulario de sorteo a la Cámara del fuero, que asigna juzgado y número de expediente (verificable por correo o en el apartado de radicaciones). La distribución de competencia es de orden público: no puede elegirse el juzgado, salvo recusación con causa (art. 17 CPCCN). \\

\textbf{¿En qué se diferencia el sistema de la Provincia de Buenos Aires (SCBA) del nacional?} &
No requiere formulario de sorteo previo: se completan directamente los campos, incluido el departamento judicial. El texto de la demanda se copia y pega en el formulario, la documentación se sube en archivos PDF separados por documento, y el envío es directo mediante firma digital registrada. \\

\midrule
\multicolumn{2}{p{0.90\textwidth}}{
\textbf{Resumen:} La unidad recorre el ciclo completo de un juicio de desalojo por cambio de destino: la entrevista con el cliente y la identificación de los distintos conflictos jurídicos en juego; la causal jurídica precisa del cambio de destino, distinta de la falta de pago o del vencimiento contractual; las herramientas previas al litigio (carta documento, mediación), no obligatorias pero recomendables; los requisitos de la demanda más allá del artículo 330, los límites del expediente electrónico y la exigencia de la firma hológrafa escaneada; la evolución del poder del abogado tras la reforma de 2015; las excepciones de falta de personería y de falta de legitimación pasiva; la distinción entre fuente y medio de prueba como factor decisivo del resultado del proceso; el planteo de inconstitucionalidad y la reserva del caso federal; y, finalmente, el trámite operativo de presentación electrónica ante el Poder Judicial de la Nación y ante la Provincia de Buenos Aires.
} \\
\bottomrule
\end{longtable}

\end{document}
